% !TeX root = ../papers/nc-p3.tex
\section{台143の二曲線の全係数とその周期一点}
\label{sec:parallel-septic-curve}
\begin{theorem}\label{thm:parallel-septic-curve}
台143の全窓と、その112型主成分内の台142への交わりは
\[
 t=\pm1,\quad y\ne0,1,\quad b=y(1-y),\quad c=1,\quad
 (A,B,C,D)=\bigl(t(1-y),1,y,t(y^2-y+1)\bigr)
\]
で尽くされる。各符号で標準Koszul左右AS正則性の必要十分条件は $y^2-y-1\ne0$ である。
指定始点の無印同型類は符号と $y$ で決まり、各符号で $y^2-y+1=0$ の二根だけが互いに同型である。
その一類は最小周期1、$y=1/2$ では2、他の正則点では4である。
二位置移動は $y\mapsto1-y$ を誘導する。
\end{theorem}
\begin{proof}[不正則点と中心式の特殊化]
前節の係数に代入すると $y^2-y+1=0$ でだけ $x_6$ が消えて台142になる。
$t=1$ の係数を使い、$y^2-y-1=0$ では $\mathbb Q[y]/(y^2-y-1)$ 上で四次関係空間を直接計算する。
各総色の64語に、三つの位置の六関係を全て配置して張る行を作る。数体の各係数を基底 $1,y$ での
有理二行行列に置き換えれば、商の四次元数は順に $(10,10,8,8)$ で総計36になる。
従って両根とも正則性を満たさない。捻りでこの障害は両符号に共通である。

他の全点には定理\ref{thm:parallel-double-zero-curve}の局所化と剰余体の論法を使う。
$\mathbb Q(y)$ で正符号の四次中心の文字0,1に三自由列・二自由列を得る。
同定理と同じ係数の四式を選び、四始点で長さ13までの中心商次元は $(1+z+z^2+z^3)^4$ の係数になる。
全ての消去に用いた単元の既約因子は2542個である。下に証明する $y\mapsto1-y$ の移動により
両方の表示で除かれる因子は66個、総次数169である。そのうち $y,y-1,y^2-y-1$ を除く63個それぞれで、
先頭係数を保つ素数の一次剰余根を取り、五次までの標準経路次元、四次中心の次元 $(3,2,0,0)$、
全四始点での中心商を検算する。既約な主成分の一般点が標準なので、そこから特殊化される全次数の
Hilbert下界がある。数体での核の持上げと中心商の上界から全ての共役根で中央四変数上の自由性と正則性が従う。

負符号への捻りは前節で全点共通に構成した。同じHilbert級数とKoszul左右AS性を保つ。
正符号では復元した四中心四次式は文字 $0,0,1,1$ である。二位置移動の文字置換も全位置で0と1の交換なので、
全長4の総色0,1を保つ。後の二式を平方すれば四式とも総色零となり、捻っても中心に残る。
よって負符号でも重み $4,4,8,8$ の多項式環上で各始点は自由階数1024になる。
\end{proof}

\begin{proof}[添字を保つ同型と始点移動]
第三長方形と同じ交代勾配を使う。位置1,3で両側の隣接部分trace二乗の固有値差の分子は、
両符号、全六文字対で互いに素である。これら二位置で得た固有射影を両方挿入して残る位置を縮約すれば、
位置0,2でも各文字対の16個の差の分子は単位イデアルを生成する。
従って全曲線で四つの枠が内在的である。
台143では始点を固定する台置換は恒等だけで、薄いトーラスの規格化 $B=1$ から $y=C$ は一意になる。
奇位置移動では台の置換は存在しない。二位置移動では各位置で0と1を交換するものだけがあり、
$1-y$ の正規係数との比は対角作用の重み行列の全左整数核で指標1となる。
台142の二根では、同じ有限検査を $\mathbb Q[y]/(y^2-y+1)$ で行うと二根を交換する同型と、一位置移動する同型をともに得る。
二つの符号の反対階数は下の通り互いに異なる。これで指定始点の同型と周期式を得る。
\end{proof}

\begin{corollary}
$u=y^2-y$ とする。正符号の反対階数は $u=-1$ で8、$u=1$ で9、他で10である。
負符号の反対階数は $u^3-u^2+3u+1=0$ で15、他で16である。
両符号で、正則な台143の各点は112型主成分全体の正則開集合に属する。
\end{corollary}
\begin{proof}
正符号の反対切断の色0,2,3は全て定階数2である。色1では四次小行列式と三次小行列式の最大公約数が
単元を除き $(u-1)(u+1)^2$ と $u+1$ で、二次小行列式は共通零点を持たない。
負符号では色0,2,3が可逆、色1の行列式分子が上記三次式で、三次小行列式の最大公約数は1である。
開近傍については正符号で、三変数族の二零条件の曲線と同様、経路基底と文字0,1の四次中心核が
一般階数を持つので局所的に四式を延長して中心商の次元上界を保てる。
負符号では112型全体の符号を交換する捻りでこの開近傍を移す。
\path{computations/verify_parallel_septic_curve.py} が二つの数体の計算、共通固有枠、台置換、
因子の局所化と全残余因子での持上げを再現する。
\end{proof}
